%% ============================================================
%% Affective Computing — Week 7 Study Notes
%% ============================================================
\documentclass[11pt,a4paper]{article}

\usepackage[a4paper, left=1.8cm, right=1.8cm, top=1.3cm, bottom=1.8cm]{geometry}
\usepackage{booktabs}
\usepackage{tabularx}
\usepackage{array}
\usepackage{enumitem}
\usepackage{titlesec}
\usepackage{fancyhdr}
\usepackage{fontenc}
\usepackage{inputenc}
\usepackage{microtype}
\usepackage{parskip}
\usepackage{hyperref}
\usepackage{amsmath}
\usepackage{amssymb}

\hypersetup{hidelinks, colorlinks=false}

%% ── Table helpers ─────────────────────────────────────────
\newcolumntype{L}[1]{>{\raggedright\arraybackslash}p{#1}}

%% ── Section headings ──────────────────────────────────────
\titleformat{\section}{\large\bfseries}{}{0em}{}[\titlerule]
\titleformat{\subsection}{\normalsize\bfseries}{}{0em}{}
\titleformat{\subsubsection}{\normalsize\bfseries}{}{0em}{}

\titlespacing{\section}{0pt}{12pt}{4pt}
\titlespacing{\subsection}{0pt}{8pt}{3pt}
\titlespacing{\subsubsection}{0pt}{6pt}{2pt}

%% ── Page style ────────────────────────────────────────────
\pagestyle{fancy}
\fancyhf{}
\renewcommand{\headrulewidth}{0pt}
\fancyfoot[C]{\small Affective Computing --- Week 7 \textbar\ Page \thepage}

%% ── Misc helpers ──────────────────────────────────────────
\newcommand{\bb}[1]{\textbf{#1}}
\setlength{\parskip}{4pt}

%% ═══════════════════════════════════════════════════════════
\begin{document}
\setlength{\arrayrulewidth}{0.4pt}

\begin{center}
\bfseries\LARGE Affective Computing — Week 7\par
\vspace{0.2cm}
\large Assignment 7 Detailed Solutions
\end{center}
\vspace{0.5cm}

%% ════════════════════════════════════════════════════════════
\section*{ASSIGNMENT 7 --- Full Solutions with Explanations}

\subsubsection*{Question 1}
\textit{Physiological signals are especially valuable in affect detection because they:}
\medskip

\begin{tabular}{L{0.3cm} L{14.9cm}}
& \textbf{A.\ Cannot be voluntarily controlled and reflect pure emotional states \checkmark} \\
& B.\ Require full attention from the participant \\
& C.\ Are easier to fake than facial expressions \\
& D.\ Are unaffected by wearable technologies \\
\end{tabular}

\vspace{0.3cm}
\noindent\textbf{Ans:} A --- Cannot be voluntarily controlled and reflect pure emotional states\\[0.2cm]
\textbf{Why this answer:} \\
Unlike facial expressions or self-reports, \textbf{physiological signals} such as heart rate, galvanic skin response (GSR), and respiration are governed by the \textbf{autonomic nervous system} and are largely \textbf{involuntary}. This makes them highly valuable for affect detection because they are difficult to consciously fake or suppress, providing a more authentic window into a person's emotional state. This involuntary nature is precisely what gives physiological signals an advantage in affective computing applications.
\vspace{0.4cm}

\subsubsection*{Question 2}
\textit{Which of the following is not a physiological measure?}
\medskip

\begin{tabular}{L{0.3cm} L{14.9cm}}
& A.\ Respiration \\
& B.\ Electrocardiogram \\
& C.\ Heart Rate \\
& \textbf{D.\ Eye-blink frequency \checkmark} \\
\end{tabular}

\vspace{0.3cm}
\noindent\textbf{Ans:} D --- Eye-blink frequency\\[0.2cm]
\textbf{Why this answer:} \\
\textbf{Respiration}, \textbf{Electrocardiogram (ECG)}, and \textbf{Heart Rate} are all direct physiological measures reflecting activity of the body's internal systems (cardiovascular, respiratory). \textbf{Eye-blink frequency}, while measurable and emotionally relevant, is a \textbf{behavioral/oculomotor measure} rather than a physiological one in the strict sense --- it is observed externally and is more susceptible to voluntary control. It falls under the category of behavioral or facial/ocular signals rather than internal physiological signals.
\vspace{0.4cm}

\subsubsection*{Question 3}
\textit{Heart rate is strongly associated with which dimension of emotion?}
\medskip

\begin{tabular}{L{0.3cm} L{14.9cm}}
& A.\ Valence \\
& \textbf{B.\ Arousal \checkmark} \\
& C.\ Dominance \\
& D.\ Surprise \\
\end{tabular}

\vspace{0.3cm}
\noindent\textbf{Ans:} B --- Arousal\\[0.2cm]
\textbf{Why this answer:} \\
Within the \textbf{dimensional model of emotion} (Russell's Circumplex Model; the VAD model), \textbf{heart rate} is most strongly associated with the \textbf{arousal dimension} --- the degree of activation or excitation in an emotional state. High-arousal emotions (fear, excitement, anger) tend to increase heart rate, while low-arousal emotions (calm, relaxation, sadness) tend to lower it. Heart rate is a poor discriminator for \textbf{valence} (positive vs.\ negative) since both fear and excitement can raise heart rate similarly.
\vspace{0.4cm}

\subsubsection*{Question 4}
\textit{A key advantage of PPG (Photo-Plethysmography) over ECG is that PPG:}
\medskip

\begin{tabular}{L{0.3cm} L{14.9cm}}
& A.\ Provides greater spatial resolution \\
& \textbf{B.\ Is quicker to attach and less intrusive \checkmark} \\
& C.\ Measures electrical pulses directly \\
& D.\ Requires medical-grade gel electrodes \\
\end{tabular}

\vspace{0.3cm}
\noindent\textbf{Ans:} B --- Is quicker to attach and less intrusive\\[0.2cm]
\textbf{Why this answer:} \\
\textbf{PPG (Photo-Plethysmography)} measures blood volume changes in the microvascular bed of tissue using light (typically via a fingertip or wrist sensor). Compared to \textbf{ECG}, which requires placement of multiple electrodes on the chest with conductive gel, PPG is \textbf{far easier to attach, less intrusive, and more comfortable} for everyday use --- making it ideal for wearable affective computing applications. The trade-off is that PPG is more susceptible to motion artifacts and is an indirect measure of cardiac activity.
\vspace{0.4cm}


\subsubsection*{Question 5}
\textit{Heart Rate Variability (HRV) measures:}
\medskip

\begin{tabular}{L{0.3cm} L{14.9cm}}
& A.\ The average heart rate over the day \\
& B.\ The electrical intensity of heart contractions \\
& \textbf{C.\ The variation in RR intervals between beats \checkmark} \\
& D.\ The number of peaks in each heartbeat waveform \\
\end{tabular}

\vspace{0.3cm}
\noindent\textbf{Ans:} C --- The variation in RR intervals between beats\\[0.2cm]
\textbf{Why this answer:} \\
\textbf{Heart Rate Variability (HRV)} refers to the \textbf{beat-to-beat variation in the time intervals between consecutive heartbeats} --- specifically the \textbf{RR intervals} (the time between two successive R-peaks in the ECG waveform). HRV is a rich indicator of autonomic nervous system activity: high HRV generally indicates good autonomic flexibility and a relaxed state, while low HRV is associated with stress, fatigue, or high arousal. HRV is widely used as a biomarker for emotional and cognitive states in affective computing.
\vspace{0.4cm}

\subsubsection*{Question 6}
\textit{A major limitation of ECG/PPG for affect recognition is that they:}
\medskip

\begin{tabular}{L{0.3cm} L{14.9cm}}
& A.\ Only measure emotional valence, not arousal \\
& B.\ Cannot measure changes over short intervals \\
& C.\ Are too invasive for most participants \\
& \textbf{D.\ Cannot distinguish positive from negative arousal \checkmark} \\
\end{tabular}

\vspace{0.3cm}
\noindent\textbf{Ans:} D --- Cannot distinguish positive from negative arousal\\[0.2cm]
\textbf{Why this answer:} \\
ECG and PPG both track \textbf{arousal-related changes} in the cardiovascular system well --- elevated heart rate signals high arousal. However, they \textbf{cannot reliably distinguish between positive high-arousal states} (e.g., excitement, joy) \textbf{and negative high-arousal states} (e.g., fear, anger), because both produce similar cardiovascular signatures. This \textbf{valence ambiguity} is a fundamental limitation, often addressed by combining cardiac signals with other modalities like GSR, facial expressions, or EEG to achieve full emotional characterization.
\vspace{0.4cm}
\newpage
\subsubsection*{Question 7}
\textit{Physiological signals like HR and GSR are less likely to be consciously manipulated compared to facial expressions.}
\medskip

\begin{tabular}{L{0.3cm} L{14.9cm}}
& \textbf{A.\ True \checkmark} \\
& B.\ False \\
\end{tabular}

\vspace{0.3cm}
\noindent\textbf{Ans:} A --- True\\[0.2cm]
\textbf{Why this answer:} \\
\textbf{Heart Rate (HR)} and \textbf{Galvanic Skin Response (GSR)} are mediated by the \textbf{autonomic nervous system} --- a system that operates largely below the threshold of conscious control. While trained individuals (e.g., through biofeedback) can exert limited influence over these signals, in general they are far \textbf{less susceptible to deliberate manipulation} than facial expressions, which involve voluntary muscles and can be consciously posed or suppressed. This is why physiological signals are considered more reliable ground-truth indicators of emotional states.
\vspace{0.4cm}

\subsubsection*{Question 8}
\textit{ECG-based heart rate measurements are more accurate than PPG because:}
\medskip

\begin{tabular}{L{0.3cm} L{14.9cm}}
& \textbf{A.\ They measure electrical activity directly \checkmark} \\
& B.\ They require fewer sensors \\
& C.\ They work better at long distances \\
& D.\ They are less sensitive to motion artifacts \\
\end{tabular}

\vspace{0.3cm}
\noindent\textbf{Ans:} A --- They measure electrical activity directly\\[0.2cm]
\textbf{Why this answer:} \\
\textbf{ECG (Electrocardiogram)} directly measures the \textbf{electrical signals} generated by the heart's depolarization and repolarization cycles, producing a precise waveform (including the characteristic P-QRS-T complex). This direct electrical measurement provides a highly accurate representation of each heartbeat. \textbf{PPG}, by contrast, indirectly infers heart rate from \textbf{optical changes in blood volume} --- making it slightly less precise and more vulnerable to motion artifacts and signal noise. ECG remains the gold standard for cardiac measurement in affective computing research.
\vspace{0.4cm}

\subsubsection*{Question 9}
\textit{GSR can reveal whether the emotional state is positive or negative.}
\medskip

\begin{tabular}{L{0.3cm} L{14.9cm}}
& A.\ True \\
& \textbf{B.\ False \checkmark} \\
\end{tabular}

\vspace{0.3cm}
\noindent\textbf{Ans:} B --- False\\[0.2cm]
\textbf{Why this answer:} \\
\textbf{GSR (Galvanic Skin Response)}, also known as \textbf{Electrodermal Activity (EDA)} or \textbf{Skin Conductance Response (SCR)}, measures changes in the electrical conductance of the skin caused by sweat gland activity --- which is modulated by sympathetic nervous system arousal. GSR is a sensitive indicator of \textbf{arousal level} (high arousal $\rightarrow$ increased GSR) but \textbf{cannot distinguish between positive and negative emotions} --- both excitement (positive) and fear (negative) can produce identical GSR elevations.
\vspace{0.4cm}
\newpage
\subsubsection*{Question 10}
\textit{Low HRV is typically associated with:}
\medskip

\begin{tabular}{L{0.3cm} L{14.9cm}}
& A.\ Relaxation \\
& \textbf{B.\ Strong stress or arousal \checkmark} \\
& C.\ Enhanced emotional control \\
& D.\ Deep sleep \\
\end{tabular}

\vspace{0.3cm}
\noindent\textbf{Ans:} B --- Strong stress or arousal\\[0.2cm]
\textbf{Why this answer:} \\
\textbf{Low Heart Rate Variability (HRV)} indicates that the time intervals between heartbeats are becoming more uniform and less variable --- a signature of \textbf{heightened sympathetic nervous system activity} (the ``fight-or-flight'' response). This is typically observed during states of \textbf{stress, anxiety, or high arousal}, where the sympathetic system dominates and overrides the naturally variable, parasympathetically-driven rhythm. Conversely, \textbf{high HRV} is associated with relaxation, emotional regulation, and parasympathetic dominance.
\vspace{0.4cm}

\medskip
\subsection*{Assignment Quick Answer Summary}

\begin{center}
\renewcommand{\arraystretch}{1.4}
\begin{tabular}{|L{0.8cm}|L{6.2cm}|L{8.2cm}|}
\hline
\bb{Q\#} & \bb{Topic} & \bb{Correct Answer} \\
\hline
1 & Value of Physiological Signals & A --- Cannot be voluntarily controlled \\
2 & Not a Physiological Measure & D --- Eye-blink frequency \\
3 & Heart Rate \& Emotion Dimension & B --- Arousal \\
4 & PPG Advantage over ECG & B --- Quicker to attach, less intrusive \\
5 & HRV Definition & C --- Variation in RR intervals \\
6 & Limitation of ECG/PPG & D --- Cannot distinguish positive/negative arousal \\
7 & HR \& GSR vs.\ Facial Expressions & True --- Less susceptible to manipulation \\
8 & ECG Accuracy over PPG & A --- Measures electrical activity directly \\
9 & GSR \& Valence & False --- GSR cannot reveal valence \\
10 & Low HRV Association & B --- Strong stress or arousal \\
\hline
\end{tabular}
\end{center}

\medskip
{\small\textit{Reference: Affective Computing --- Week 7 Lecture Materials}}

\end{document}
